\section*{第 \S\arabic{exercisegroup} 节习题}
\phantomsection
\addcontentsline{toc}{section}{第 \protect\S\arabic{exercisegroup} 节习题}

\begin{enumerate}
\def\labelenumi{\arabic{enumi})}
\item
  定义一个递增函数 \(g\) ，它在 \(+\infty\) 的某个邻域上有连续导数，使 \(g\) 与 \(1 / x\) 可比较到 1 阶，但 \(x\) 与 \(1 / g\) 不可比较到 1 阶（取 \(g'(x) = 1\) ，但在以点 \(x = n\) （ \(n\) 为整数 \(>0\) ）为中心的充分小区间上 \(g'\) 取非常大的值）。
\item
  设 f 与 g 是两个 \(\geqslant0\) 的函数（随 x 趋于 \(+\infty\) ），且可比较到 1 阶；若 h 是一个可导函数， \(\geqslant0\) 且递增（当 x 趋于 \(+\infty\) ），证明 hf 与 hg 可比较到 1 阶。
\item
  给出两个函数的例子：它们为正、递减、当 x 趋于 \(+\infty\) 时趋于 0、可比较到 1 阶，且使 \(xf(x)\) 与 \(xg(x)\) 不可比较到 1 阶（取 f 与 g 等价，且使
\end{enumerate}

\[
\frac {f ^ {\prime} (x)}{f (x)} + \frac {1}{x} \qquad \text { and } \qquad \frac {g ^ {\prime} (x)}{g (x)} + \frac {1}{x}
\]

不可比较）。

\begin{enumerate}
\def\labelenumi{\arabic{enumi})}
\setcounter{enumi}{3}
\tightlist
\item
  \begin{enumerate}
  \def\labelenumii{\alph{enumii})}
  \tightlist
  \item
    证明 \(\frac{\sin x}{\sqrt{x}} \sim \frac{\sin x}{\sqrt{x}} \left(1 + \frac{\sin x}{\sqrt{x}}\right)\) ，但积分 \(\int_{a}^{+\infty} \frac{\sin t}{\sqrt{t}} dt\) 收敛，而积分 \(\int_{a}^{+\infty} \frac{\sin t}{\sqrt{t}} \left(1 + \frac{\sin t}{\sqrt{t}}\right) dt\) 为无穷。
  \end{enumerate}
\end{enumerate}

\begin{enumerate}
\def\labelenumi{\alph{enumi})}
\setcounter{enumi}{1}
\item
  证明积分 \(\int_{a}^{+\infty}\frac{\sin t}{t}dt\) 与 \(\int_{a}^{+\infty}\frac{\sin^{2}t}{t^{2}}dt\) 都收敛，但 \(\int_{x}^{+\infty}\frac{\sin^{2}t}{t^{2}}dt\) 相对于 \(\int_{x}^{+\infty}\frac{\sin t}{t}dt\) 并非可忽略，尽管 \(\sin^{2}t/t^{2}\ll\sin t/t\) 。
\item
  考虑取值于 \(R^{2}\) 的两个连续函数（定义在 \([1, +\infty[\) 上）：
\end{enumerate}

\[
\mathbf {f}: x \mapsto \left(\frac {1}{x ^ {2}}, \frac {\sin x}{x}\right), \quad \mathbf {g}: x \mapsto \left(\frac {1}{x ^ {2}}, \frac {\sin x}{x} \left(1 + \frac {\sin x}{x}\right)\right).
\]

证明它们对 \(x\) 的任何值都不为零， \(\mathbf{f} \sim \mathbf{g}\) （当 \(x\) 趋于 \(+\infty\) 时），积分 \(\int_{a}^{+\infty} \mathbf{f}(t) dt\) 与 \(\int_{a}^{+\infty} \mathbf{g}(t) dt\) 收敛，但 \(\int_{x}^{+\infty} \mathbf{f}(t) dt\) 不等价于 \(\int_{x}^{+\infty} \mathbf{g}(t) dt\) （当 \(x\) 趋于 \(+\infty\) 时）。

5) 设 \(f\) 是定义在 \(+\infty\) 的某个邻域上的实凸函数，且满足 \(f(x) \gg x\) 。称 \(f\) 在 \(+\infty\) 的某个邻域上是正则凸的，如果对每个凸函数 \(g\) （定义在 \(+\infty\) 的某个邻域上）且满足 \(f \sim g\) ，也有 \(f_d' \sim g_r'\) 。

对每个数 \(\alpha > 0\) 与充分大的 x，设 \(k(\alpha, x)\) 是诸数 \(\left(f(y) - (y - x)f_{d}'(y)\right)/f(x)\) 的下确界（当 y 取遍 \(\geqslant x\) 且满足 \(f_{d}'(y) \leqslant (1 + \alpha)f_{d}'(x)\) 的数所成的集时）。类似地，设 \(h(\alpha, x)\) 是诸数

\[
\left(f (z) + (x - z) f _ {d} ^ {\prime} (z)\right) / f (x)
\]

（当 \(z\) 取遍 \(\leqslant x\) 且满足 \(f_d'(z) \geqslant (1 - \alpha)f_d'(x)\) 的数所成的集时）的下确界。设 \(\psi(\alpha) = \limsup_{x \to +\infty} k(\alpha, x)\) ， \(\varphi(\alpha) = \limsup_{x \to +\infty} h(\alpha, x)\) 。证明： \(f\) 在

\(+\infty\) 的某个邻域上正则凸的充分必要条件是，对一切充分小的 \(\alpha\) 有 \(\psi(\alpha) < 1\) 且 \(\varphi(\alpha) < 1\) 。

6) 设 \(\mathbf{f}\) 是某个区间 \([x_0, +\infty[\) （属于 \(\mathbf{R}\) ）上的连续向量函数，且对所有 \(\lambda \geqslant 0\) ，函数 \(\mathbf{f}(x + \lambda) - \mathbf{f}(x)\) 当 \(x\) 趋于 \(+\infty\) 时趋于 0。

\begin{enumerate}
\def\labelenumi{\alph{enumi})}
\tightlist
\item
  证明： \(\mathbf{f}(x+\lambda)-\mathbf{f}(x)\) 当 \(\lambda\) 属于任一紧区间 \(K=[a,b]\) （ \([0,+\infty[\) 中的）时随 1/x 一致地趋于 0。（用反证法：若存在一个序列 \((x_{n})\) 趋于 \(+\infty\) ，以及 K 中的点列 \((\lambda_{n})\) 使
\end{enumerate}

\[
\| \mathbf {f} (x _ {n} + \lambda_ {n}) - \mathbf {f} (x _ {n}) \| > \alpha > 0
\]

对每个 n 成立，则存在 \(J_{n}\) （ \(\lambda_{n}\) 在 K 中的一个邻域），使对每个 \(\lambda \in J_{n}\) 有 \(\|\mathbf{f}(x_{n} + \lambda) - \mathbf{f}(x_{n})\| > \alpha\) 。用归纳法构造一个递减的闭区间序列 \(I_{k} \subset K\) 与一个子序列 \((x_{n_{k}})\) （ \((x_{n})\) 的），使 \(\|\mathbf{f}(x_{n_{k}} + \lambda) - \mathbf{f}(x_{n_{k}})\| \geqslant \alpha/3\) 对每个 \(\lambda \in I_{k}\) 成立；为此注意，若 \(\delta_{k}\) 是 \(I_{k}\) 的长度而 q 是使 \(q\delta_{k} > b - a\) 的整数，则只要 x 充分大就有 \(\|\mathbf{f}(x + \delta_{k}) - \mathbf{f}(x)\| \leqslant \alpha/3q\) 。）

\begin{enumerate}
\def\labelenumi{\alph{enumi})}
\setcounter{enumi}{1}
\tightlist
\item
  由 \(a\) ) 推出 \(\int_{x}^{x + 1}\mathbf{f}(t)dt - \mathbf{f}(x)\) 当 \(x\) 趋于 \(+\infty\) 时趋于 0，并断定 \(\mathbf{f}(x) = o(x)\) 。
\end{enumerate}

\begin{enumerate}
\def\labelenumi{\arabic{enumi})}
\setcounter{enumi}{6}
\tightlist
\item
  设 g 是一个严格正的实函数，在某个区间 \([x_{0}, +\infty[\) 上连续，且对每个 \(\mu > 0\) 有 \(g(\mu x) \sim g(x)\) 。由习题 6 推出 g 相对于 x 的阶为 0。
\end{enumerate}

\setcounter{exercisegroup}{3}
\refstepcounter{exercisegroup}
